\section{Local first variation and height estimates}\label{sec:local}

We first collect the local estimates used in the signed comparison.  Their
constants depend on the dimensions, an upper bound for the integer under
consideration, fixed interior enlargement factors, and the indicated bounds
for the coordinates.  Counts in a normal fiber always include integer
multiplicity.  A configuration of size $Q$ is therefore an unordered list of
$Q$ vectors, with repetitions allowed.

\subsection{First variation and integer counts}

We begin with the standard stationary monotonicity and mean-value
arguments; compare \cite[Theorem~9]{Allard1987} for monotonicity.
We include the derivation to retain the lower mass bound at every point
of the support.

\begin{lemma}[Monotonicity and mean value]\label{lt:monotonicity}
Let $V$ be stationary and integral in an open Euclidean set.  If
$p\in\spt\mu$ and $\overline{\mathbf B_r(p)}$ lies in that set, then
\begin{equation}\label{eq:lt-lower-mass}
 \mu(\mathbf B_r(p))\ge\omega_m r^m.
\end{equation}
If $u\ge0$ is a $C^2$ ambient function with
$\tr(P_TD^2u)\ge0$ for $V$-almost every $(x,T)$ in the ball, then
\begin{equation}\label{eq:lt-mean-value}
 u(p)\le \frac{1}{\omega_m r^m}
                  \int_{\mathbf B_r(p)}u\,\dd\mu.
\end{equation}
The same assertion holds for continuous functions satisfying the corresponding
weak subharmonicity inequality.
\end{lemma}

\begin{proof}
Put $t=|x-p|$ and let $f\ge0$ be a smooth decreasing radial cutoff.  Testing
stationarity with $(x-p)f(t/r)$ and differentiating under the integral gives
\[
 \frac{\dd}{\dd\log r}\left(r^{-m}\int f(t/r)\,\dd\mu\right)
 =r^{-m}\int \frac{-f'(t/r)}{rt}
                |(I-P_T)(x-p)|^2\,\dd\mu\ge0.
\]
Approximation of the indicator of $[0,1)$ proves ball monotonicity.  At a point
where the rectifiable density is an integer, its limit at zero is at least
$\omega_m$.  Such points are dense in the support.  Applying the conclusion
at points tending to $p$, with radii increasing to $r$, proves
\eqref{eq:lt-lower-mass} at every support point.

Stationarity with $\varphi\nabla u$ shows
$\int\nabla_Tu\cdot\nabla_T\varphi\,\dd\mu\le0$ for every nonnegative
compactly supported $\varphi$.  The preceding radial calculation, with its
field multiplied by $u$, has the additional term
\[
 r^{-m}\int f(t/r)\nabla_Tu\cdot(x-p)\,\dd\mu.
\]
It is nonnegative: use the nonnegative radial primitive
$\varphi(t)=\int_t^\infty a f(a/r)\,\dd a$ in the weak inequality.
Thus $r^{-m}\int_{\mathbf B_r(p)}u\,\dd\mu$ is nondecreasing.
At integer-density points its limit is at least $\omega_m u(p)$.
Continuity and the same approximation of the center prove
\eqref{eq:lt-mean-value}.  This argument uses only the weak inequality, and
therefore proves the last assertion as well.
\end{proof}

We use a reference disk $S$ of class $C^J$, where the fixed finite
order $J\ge6$ is chosen sufficiently large for all derivatives below,
including the dimension-dependent test order in \Cref{lt:potential}.
The disk has a fixed normal tube, small graph slope, and bounded coordinate
derivatives through order $J$.  All pulled-back first-variation fields
have at least class $C^1$; Gaussian smoothing takes place in the base
variables.  The same fixed derivative order suffices under rescaling and
throughout the induction.  The flat case permits an arbitrary
vertical translate.  Restrict the varifold to a cylinder whose support is
separated from its vertical boundary over every compact subset of the base.
All fields used below have compact base support and are multiplied by a
vertical cutoff equal to one on that support.  Consequently this restriction
has no vertical boundary term.  Fixed enlargements of a working ball are
always assumed to remain in this cylinder.

Write a point of the tube as $q=y+z_{\rm phys}$, where $y=\pi(q)\in S$ and
$z_{\rm phys}\in N_yS$.  The second fundamental form is denoted by
$\mathrm{II}$, its trace by $H_S$, and its shape operators by
\[
 \langle A_v a,b\rangle=\langle\mathrm{II}(a,b),v\rangle,
 \qquad J_z=I-A_{z_{\rm phys}}.
\]
All norms of $H_S$ and $\nabla^\perp H_S$ below are in the physical Euclidean
metric.  At $q$, decompose the tangent projection $P=P_{T_qV}$ relative to
$T_yS\oplus N_yS$, and set
\begin{equation}\label{eq:lt-tensors}
 B=I-P_{hh},\qquad e=\tr B=\tr P_{vv},\qquad
 \lambda=|z_{\rm phys}|^2+|H_S|^2+|\nabla^\perp H_S|^2.
\end{equation}
Thus $B$ is a tensor at $q$, not a measure on the base.  It is positive
semidefinite, $0\le e\le m$, and the projection identities imply
\begin{equation}\label{eq:lt-projection-identities}
 P_{hv}P_{vh}=B-B^2,\qquad
 |P_{vh}a|^2\le B[a,a],\qquad |P_{vh}|\le C\sqrt e.
\end{equation}
In physical coordinates the measures on the base are
\begin{equation}\label{eq:lt-measures}
 \nu=\pi_\#\mu,\quad M=\nu-Q\vol_S,\quad
 D=\pi_\#(e\mu),\quad \Lambda=\pi_\#(\lambda\mu),\quad
 L=2M-D,\quad \tau=|M|+D+\Lambda.
\end{equation}
For a flat reference we set $\lambda=0$ and $\Lambda=0$.  In particular a
constant sheet offset is not charged as a geometric error in that case.

At the center of the original physical chart, normalize the base
coordinates so that their metric is Euclidean there.  For the coordinate
immersion $F$, let $g=(DF)^\top DF$, choose the orthonormal horizontal
frame $E=DF\,g^{-1/2}$ using the positive square root, and write
$\mathsf C_y=g^{1/2}$; thus $DF\,X=E\mathsf C_yX$.
Choose the normal frame by parallel transport along radial geodesics
from the same center.  This normal frame is the radial gauge for the
original chart; its regularity and connection bounds follow from the
fixed geometry above.
When a physical base radius $l$ is scaled to one, we use base coordinates
$y/l$ and height coordinates $\zeta=z_{\rm phys}/l$ in this fixed normal
frame.  The fixed coordinate bounds and the normalization at the center give
\[
 |\mathsf C_y-I|+|D_{y/l}\mathsf C_y|+|D_{y/l}E|\le Cl
\]
throughout the original chart.  Measures, including the reference volume,
are divided by $l^m$ and pushed forward to these coordinates.  We retain the
symbols in
\eqref{eq:lt-measures} for the resulting measures.  The function $\lambda$
is still evaluated in physical units; it is not replaced by
$|\zeta|^2$ or by rescaled curvature.  A further coordinate ball of radius
$s$ consequently has physical radius $ls$.  Define
\begin{equation}\label{eq:lt-small-geometry}
 \gamma=l+\sup|z_{\rm phys}|.
\end{equation}
The reference disks used below have a uniformly small tube, so that $J_z$
and $J_z^{-1}$ are uniformly bounded.

\begin{lemma}[Curved horizontal first variation]\label{lt:curvedvariation}
For a compactly supported tangent field $X$ on $S$,
\begin{equation}\label{eq:lt-curvedvariation}
 \begin{split}
 &\left|\int\Div_SX\,\dd M-\int B:\nabla^SX\,\dd\mu\right|\\
 &\quad\le C\int\Big\{|z_{\rm phys}|e|\nabla^SX|
 +\big[(|\nabla^\perp H_S|+|z_{\rm phys}|)|z_{\rm phys}|
       +\sqrt e\,|z_{\rm phys}|\big]|X|\Big\}\dd\mu.
 \end{split}
\end{equation}
The integrands containing tensors are evaluated at $q$, with the base field
pulled back by $\pi$.  In coordinates scaled by $l$, the right side is
bounded by
\begin{equation}\label{eq:lt-scaledvariation}
 C\gamma\int(e+\lambda)(|\nabla X|+|X|)\,\dd\mu.
\end{equation}
Here and below a dimensionless scaled test field is used in a scaled
identity.
\end{lemma}

\begin{proof}
The nearest-point derivative is $D\pi=J_z^{-1}P_h$.  Lift $X$ to the ambient
field $\overline X(q)=J_zX(y)$.  This is the velocity obtained by moving $y$
with velocity $X$ and transporting $z_{\rm phys}$ parallel in the normal
connection.  Compute at a point in tangent and normal frames whose respective
connections vanish there.  With
$K_z(a)=(\nabla_a A)_{z_{\rm phys}}X$, the four blocks of its derivative are
\begin{align*}
 (D\overline X)_{hh}&=J_z(\nabla X)J_z^{-1}-K_zJ_z^{-1},&
 (D\overline X)_{hv}v&=-A_vX,\\
 (D\overline X)_{vh}a&=\mathrm{II}(J_z^{-1}a,J_zX),&
 (D\overline X)_{vv}&=0.
\end{align*}
The first term of the horizontal block has trace $\Div_SX$.  Codazzi in
Euclidean space gives
\[
 \tr K_z=\langle\nabla_X^\perp H_S,z_{\rm phys}\rangle.
\]
For completeness, Codazzi follows by writing
$\mathrm{II}_{ij}=(\partial_i\partial_jF)^\perp$ in local coordinates and
commuting the three derivatives of the immersion; at the chosen point the
connection terms are precisely those in the covariant derivative.
Since $K_z=O(|z_{\rm phys}||X|)$ and $J_z^{-1}-I=O(|z_{\rm phys}|)$,
\[
 \tr(D\overline X)_{hh}
 =\Div_SX-\langle\nabla_X^\perp H_S,z_{\rm phys}\rangle
       +O(|z_{\rm phys}|^2|X|).
\]
The symmetric sum of the off-diagonal blocks, paired with $a\in T_yS$ and
$v\in N_yS$, is exactly
\[
 \langle a,J_z^{-1}[A_{z_{\rm phys}},A_v]X\rangle.
\]
Indeed it equals
$\langle a,(J_z^{-1}A_vJ_z-A_v)X\rangle$.
Its size is $O(|z_{\rm phys}||a||v||X|)$, including in arbitrary codimension.
Contracting it with $P_{vh}$ costs
$C\sqrt e\,|z_{\rm phys}||X|$ by
\eqref{eq:lt-projection-identities}.  Contracting the horizontal block with
$I-B$ instead of the identity gives $-B:\nabla X$ and an error at most
$C|z_{\rm phys}|e(|\nabla X|+|X|)$.  Absorb its last summand in the
preceding $\sqrt e$ term.  Integrating $P:D\overline X=0$ and using
$\int_S\Div_SX\,\dd\vol_S=0$ proves
\eqref{eq:lt-curvedvariation}.

After scaling, the coefficients of $|\nabla X|$ and $|X|$ are respectively
\[
 |z_{\rm phys}|e,\qquad
 l\big[(|\nabla^\perp H_S|+|z_{\rm phys}|)|z_{\rm phys}|
                         +\sqrt e\,|z_{\rm phys}|\big].
\]
Young's inequality bounds them by $C\gamma(e+\lambda)$.  Bounded coordinate
changes preserve this estimate.
\end{proof}

\begin{lemma}[Projection and integer counts]\label{lt:jacobian}
The tangential Jacobian of $\pi$ is
\begin{equation}\label{eq:lt-jacobian}
 J_\pi=\frac{\sqrt{\det P_{hh}}}{\det J_z},\qquad
 |J_\pi-1|\le C(e+\lambda).
\end{equation}
There is a nonnegative integer-valued function $N$ on the base, finite
almost everywhere, such that
\begin{equation}\label{eq:lt-integer-measure}
 \pi_\#(J_\pi\mu)=N\vol_S,
 \qquad |\nu-N\vol_S|\le C(D+\Lambda).
\end{equation}
The inequalities are measure inequalities on every relatively compact base
set.  Moreover, for sufficiently small $e$,
\begin{equation}\label{eq:lt-signed-jacobian}
 2(1-J_\pi)-e\le Ce^2+C(|z_{\rm phys}|^2+|H_S|^2).
\end{equation}
\end{lemma}

\begin{proof}
The squared Jacobian of the orthogonal projection of $T$ onto $T_yS$ is
$\det P_{hh}$.  Composition with $J_z^{-1}$ proves the formula.  If the
eigenvalues of $B$ are $b_i\in[0,1]$, then
$1-\sqrt{\prod_i(1-b_i)}\le\sum_i b_i=e$.
Also
\[
 \det J_z=1-\langle H_S,z_{\rm phys}\rangle+O(|z_{\rm phys}|^2).
\]
These facts and Young's inequality prove the bound in
\eqref{eq:lt-jacobian}.  Expanding the square root one order further gives
$\sqrt{\det(I-B)}=1-e/2+O(e^2)$, which proves
\eqref{eq:lt-signed-jacobian}.  The rectifiable area formula
\cite[Chapter~3, equation~(2.5)]{Simon2018} gives the first equality in
\eqref{eq:lt-integer-measure}, where $N$ counts points of nonzero projection
Jacobian with their multiplicities.  Its second assertion follows by pushing
forward $(1-J_\pi)\mu$.  The part on which the projection has deficient rank
has $e\ge1$ and is included in this error; no absolute-continuity assertion
about its pushforward is needed.
\end{proof}

The next elementary analytic observation explains why a scalar
weak-$L^1$ estimate is sufficient here.  Both its distributional estimate and
the integer count are essential. Its proof uses the Calder\'on--Zygmund
decomposition and cancellation argument
\cite[Chapter~I, Lemmas~1--2]{CalderonZygmund1952}, with a localized
potential and a harmonic remainder controlled modulo constants.

\begin{lemma}[Scalar potential estimate]\label{lt:potential}
Suppose a finite scalar measure $v$ on $B_4\subset\R^m$ satisfies
\begin{equation}\label{eq:lt-gradient-equation}
 \partial_jv=\sum_i\partial_i F_{ij}+g_j,
 \qquad \sum_{ij}|F_{ij}|(B_4)+\sum_j|g_j|(B_4)\le\epsilon.
\end{equation}
Then there is a real constant $c$ such that the density $v_{\rm ac}$ obeys
\begin{equation}\label{eq:lt-weak-one}
 |\{y\in B_2:|v_{\rm ac}(y)-c|>t\}|
                     \le C\epsilon/t\quad(t>0),
\end{equation}
and, for every smooth test supported in $B_2$,
\begin{equation}\label{eq:lt-distribution-control}
 |\langle v-c\,\dd y,\varphi\rangle|
                     \le C\epsilon\|\varphi\|_{C^{k}},
\end{equation}
where $k$ is a fixed integer depending only on $m$.  Fixed changes of the
radii give the same conclusion.
\end{lemma}

\begin{proof}
First mollify on a slightly smaller domain.  Choose a cutoff equal to one on
$B_3$ and supported in $B_{7/2}$, and let $\Phi$ be the fundamental solution
of the Laplacian.  The potential
\[
 u=\sum_{ij}\partial_i\partial_j\Phi*(\chi F_{ij})
                       +\sum_j\partial_j\Phi*(\chi g_j)
\]
has the same Laplacian as $v$ on $B_3$, so $h=v-u$ is harmonic there.
The second-derivative kernels have the weak-type estimate
$|\{|u|>t\}\cap B_3|\le C\epsilon/t$; the first-derivative kernels are
integrable on bounded sets.  We recall the kernel argument.  Decompose a
mollified finite measure at level $t$ into its bounded part and zero-mean
pieces on maximal dyadic cubes.  The bounded part has squared $L^2$ norm at
most $Ct$ times the original mass norm.  The Fourier multiplier
$\xi_i\xi_j/|\xi|^2$ has $L^2$ norm at most one.  The bad cubes have total
volume at most $C\epsilon/t$; off their fixed enlargements, cancellation and
the bound $|\nabla\partial_i\partial_j\Phi(x)|\le C|x|^{-m-1}$ bound the
$L^1$ norm of their contributions by $C\epsilon$.  Chebyshev's inequality
proves the stated weak estimate, uniformly in mollification.  Pairing the
same potentials against a smooth test, moving the derivatives to the test,
also gives $|\langle u,\varphi\rangle|\le
C\epsilon\|\varphi\|_{C^k}$ on any fixed bounded region.

The harmonic remainder is controlled modulo a constant by the full gradient
equation.  For $x\in B_{5/2}$ choose a fixed radial unit-integral bump
$\eta_x$ supported in a small ball about $x$.  Harmonic mean value and
\eqref{eq:lt-gradient-equation} give
\[
 \partial_jh(x)=\langle\partial_jh,\eta_x\rangle
   =-\sum_i\langle F_{ij},\partial_i\eta_x\rangle
      +\langle g_j,\eta_x\rangle
      +\langle u,\partial_j\eta_x\rangle.
\]
Consequently $\sup_{B_{5/2}}|\nabla h|\le C\epsilon$.  Taking $c=h(0)$
proves both estimates for the mollified quantities.  The constants $c$ are
bounded by pairing against one fixed unit-integral bump and the finite mass
of $v$.  Pass to a subsequence as the mollification tends to zero; the
densities converge almost everywhere at Lebesgue points.  This proves
\eqref{eq:lt-weak-one}, and weak convergence against tests proves
\eqref{eq:lt-distribution-control}.
\end{proof}

\begin{lemma}[Scalar mass and propagation of the integer]\label{lt:scalar}
Fix a bound for $Q$ and fixed concentric interior enlargements of a ball
$B_s$.  Suppose the normalized projected mass on the outer enlargement is
bounded, and
\[
 \epsilon=s^{-m}(D+\Lambda)(B_{cs})
\]
is sufficiently small; $c>1$ is a fixed enlargement with room for the
argument.  The count selects a unique integer $q\ge0$ on most of an
intermediate ball.  If either a fixed smooth mass average differs from $Q$
by less than $1/4$, or $N=Q$ on a fixed majority of that ball, then $q=Q$.
In that case
\begin{equation}\label{eq:lt-scalar-mass}
 |M|(B_s)+\vol_S\{y\in B_s:N(y)\ne Q\}
                       \le C(D+\Lambda)(B_{cs}).
\end{equation}
In the smooth-average condition the test is nonnegative, has reference
integral one, and has fixed scaled derivative bounds.  The same integer
propagates along overlapping comparable balls for which these small-energy
and mass bounds hold and whose overlap has a fixed positive relative
volume.
\end{lemma}

\begin{proof}
Scale to $s=1$.  Write $\dd\vol_S=a(y)\,\dd y$ and put the factor $a^{-1}$
in each coordinate component of the tangent test field.  Then
\Cref{lt:curvedvariation} gives an equation of the form
\eqref{eq:lt-gradient-equation} for the scalar measure $a^{-1}\nu$.
The divergence data have mass at most $C(D+\Lambda)$, as do the
zeroth-order data; derivatives of the fixed coordinate factors have bounded
size.  In particular, no derivative of the unknown scalar mass is estimated
separately.  Apply \Cref{lt:potential}.  By \Cref{lt:jacobian}, the difference
between $a^{-1}\nu$ and $N\,\dd y$ has mass at most $C\epsilon$.
The density version of this difference and \eqref{eq:lt-weak-one} give
\[
 |\{|N-c|>t\}\cap B_2|\le C\epsilon/t.
\]
If $\epsilon=0$, this estimate for every $t>0$ gives $N=c$ almost
everywhere.  Since $N$ is integer-valued, $c$ is an integer, and the
distributional bound is exact; the conclusions follow directly.
Assume henceforth that $\epsilon>0$.
Taking $t=C_0\epsilon$ with $C_0$ large proves that, on a fixed large
majority, $N$ is within $C_0\epsilon$ of $c$.  For small $\epsilon$ there
is at most one integer in this interval; call it $q$.  Thus
$|c-q|\le C\epsilon$.  The distributional bound against the specified
smooth average makes its value $q+O(\epsilon)$, so the smooth-average
hypothesis forces $q=Q$.  The majority hypothesis has the same consequence
by intersecting the two majority sets.

Take now $t$ to be a fixed fraction of one.  Since $|c-Q|\le C\epsilon$,
\begin{equation}\label{eq:lt-count-exception}
 |\{N\ne Q\}\cap B_2|\le C\epsilon.
\end{equation}
The nonnegativity of $N$ gives
$\int_{B_2}(N-Q)_-\,\dd y\le Q|\{N\ne Q\}\cap B_2|$.
The small measure error between $\nu$ and $N\vol_S$ therefore bounds
$M_-(B_2)$ by $C\epsilon$.  Choose $0\le\varphi\le1$, equal to one on
$B_1$ and supported in $B_2$.  Equation
\eqref{eq:lt-distribution-control}, applied also to the bounded volume
factor, gives $|\int\varphi\,\dd M|\le C\epsilon$.  Hence
\[
 M_+(B_1)\le\int\varphi\,\dd M+\int\varphi\,\dd M_-
                                                    \le C\epsilon.
\]
This proves total variation and the exceptional-volume assertion, and
rescaling proves \eqref{eq:lt-scalar-mass}.  Finally, on an overlap of fixed
relative volume, the same fixed-threshold argument gives
$|\{N\ne q_i\}|\le C\epsilon_i$ for each ball's selected integer $q_i$,
before any identification with $Q$.  These exceptional sets cannot fill
the overlap.  At a remaining point both selected integers equal $N$.
Induction along a finite chain proves propagation.
\end{proof}

\begin{remark}\label{rem:lt-scalar-radii}
The enlargement constant in \eqref{eq:lt-scalar-mass} can be any prescribed
$c>1$, at the cost of constants depending on the interior room.  Applying
the argument on a finite cover gives, in particular, the notation
$|M|(B)\le C(D+\Lambda)(2B)$.  It is never an estimate on a ball with no
interior enlargement or with an unspecified integer.
\end{remark}

\subsection{Configurations and height bounds}

\begin{lemma}[Configurations, higher moments, and support heights]
\label{lt:config}
Fix $Q\ge1$, and choose a sufficiently large fixed enlargement $c>1$
depending only on the dimensions, $Q$, and the fixed coordinate and
cylinder-room data.  For every prescribed exceptional proportion
$\delta_0>0$ there is a threshold $d_0=d_0(\delta_0)>0$, depending otherwise
only on these fixed data, such that the following holds.  Suppose $Q$ is
selected on the enlarged ball $B_{cs}(x)$ and
\[
 s^{-m}\tau(B_{cs}(x))\le d\le d_0,
\]
with the cylinder and interior-room conventions above.  There is a
configuration $\mathcal A=\{a_1,\ldots,a_Q\}$ of
scaled heights such that, outside a subset of $B_s(x)$ of relative volume
at most $\delta_0$, the fiber consists of exactly $Q$ projection entries and
can be matched to $\mathcal A$ within $C_{\delta_0}s\sqrt d$.
Configurations on comparable overlapping balls can be chosen, and any such
majority configurations can be matched, within the sum of these bounds,
provided the exceptional proportions are smaller than a fixed fraction of
the overlap.

For every fixed finite $p>2$ satisfying
\begin{equation}\label{eq:lt-p-range}
 m(1/2-1/p)<1,
\end{equation}
and every fixed $0<\beta<\min\{1/2,1/m\}$, this same configuration satisfies
\begin{align}
 \left(s^{-m}\int_{\pi^{-1}(B_s(x))}
       \dist(\zeta,\mathcal A)^p\,\dd\mu\right)^{1/p}
       &\le C_p s\sqrt d,\label{eq:lt-config-moment}\\
 \sup_{\spt\mu\cap\pi^{-1}(B_s(x))}
       \dist(\zeta,\mathcal A)&\le C_\beta s d^\beta.
       \label{eq:lt-config-sup}
\end{align}
The enlargement $c$, threshold $d_0$, and configuration do not depend on
$p$ or $\beta$.  The constants $C_p,C_\beta$ may depend on their indicated
exponents and the fixed data, but not on $\delta_0$; the majority constant
$C_{\delta_0}$ may depend on it.  The support estimate holds for all support
points, including points with zero projection Jacobian.
\end{lemma}

\begin{proof}
It is enough to prove the estimates for $d>0$.  If the actual energy is
zero, fix $0<\delta'<\min\{\delta_0,1/2\}$ and positive upper bounds
$d_j\downarrow0$ below $d_0(\delta')$.  Let $G_j$ be the measurable good
sets from the positive case.  Fatou's lemma applied to their complements
gives
\[
 \vol_S(\limsup_jG_j)\ge(1-\delta')\vol_S(B_s(x)).
\]
Choose a defined fiber $\mathcal F_0$ in this limsup and retain the indices
$j_k$ for which it is good.  Its $Q$ entries match the corresponding
configurations within $C_{\delta'}s\sqrt{d_{j_k}}\to0$, so those configurations
converge to $\mathcal F_0$ up to permutations.  The limsup of the selected
$G_{j_k}$ has the same volume bound.  On each fiber in that limsup, the
matching triangle inequality along its good subsubsequence gives exact
matching to $\mathcal F_0$.  For one admissible $\beta_0>0$, the support
estimate and the matching convergence put every support height in
$\mathcal F_0$.  All support distances and hence all finite moments are
zero, proving the assertion for $d=0$.
We give the finite selection argument before the iteration that controls
exceptional fibers.  If $f$ is a bounded Lipschitz function of height with
$|f|\le1$, use $f(\zeta)\overline X$ in the horizontal first variation.
In flat coordinates the additional term is
$\nabla f\cdot P_{vh}X$.  In normal coordinates its additional frame term is
bounded by $C\Lip(f)l|z_{\rm phys}||X|$.  In the curved case, as in
\Cref{lt:scalar}, write $\dd\vol_S=a\,\dd y$ and take the coordinate
components $X=a^{-1}Y$.  Then $\Div_SX=a^{-1}\Div Y$ exactly; derivatives
of $a^{-1}$ remain in the tensor and geometric terms.  Cauchy--Schwarz,
the normalized mass bound, and \Cref{lt:potential,lt:jacobian} show that the
fiber function
\[
 N_f(y)=\sum_{q\in\pi^{-1}(y),\ J_\pi(q)>0}\theta(q)f(\zeta(q))
\]
differs from its fixed smooth average $\overline N_f$ by a weak-$L^1$
quantity whose normalized size is at most
\begin{equation}\label{eq:lt-multiplier-oscillation}
 C\{d+s\Lip(f)\sqrt d\}.
\end{equation}
The same bound controls distributional pairings.  Changing between
Jacobian-weighted and ordinary mass in this statement costs $Cd$.

For this finite selection prescribe $0<\delta\le1/4$; a larger requested
exceptional proportion is covered by using $\delta=1/4$.
Normalize the smoothly averaged vertical counting measure to a probability
measure $\mathcal P$; its total before normalization is $Q+O(d)$ by
\Cref{lt:scalar}.  Fix a mass mesh $\eta_0$ much smaller than $1/Q$.
For each of the $n$ coordinates choose its quantiles at levels
$j\eta_0$, and test smoothed half-lines immediately to both sides of these
quantiles, with transition width
$w=K s\sqrt d$.  Repeated quantiles cause no problem: retain both one-sided
tests.  There are at most $C(n)/\eta_0$ tests.  For each, the right side of
\eqref{eq:lt-multiplier-oscillation} is $C(d+K^{-1})$.
First choose $\eta_0$, then choose $K$ large enough for these tests and the
joint-count tests below, and finally choose a positive energy threshold small.
The union of the exceptional sets where any test differs from its average
by more than $\eta_0$ then has relative volume at most $\delta/3$.

Also discard the set where $N\ne Q$.  By \eqref{eq:lt-count-exception}
its relative volume is $O(d)$, so it is at most $\delta/3$ after decreasing
the same threshold.  Every remaining fiber therefore has exactly $Q$
projection entries, counted with multiplicity.

Let $F_y$ be the empirical coordinate distribution on a remaining fiber
(normalized by $Q$), and let $F$ be that of $\mathcal P$.  Bracketing an
arbitrary real number by adjacent quantiles gives
\[
 F(t-2w)-3\eta_0\le F_y(t)\le F(t+2w)+3\eta_0.
\]
If $t$ is in a quantile jump, the left and right tests at that quantile give
the same inequalities; thus atoms require no continuity assumption.
Comparison of two good fibers gives this estimate with $4w$ and
$6\eta_0$.  For $6\eta_0<1/(2Q)$, comparison at the levels
$(j-1/2)/Q$ shows that their ordered coordinate lists differ by at most
$Cw$.  Choose one such fiber as a temporary reference.  In each coordinate
group its at most $Q$ levels by joining consecutive levels whose gap is at
most a sufficiently large fixed multiple of $w$.  Each resulting interval
has length at most $CQw$, and distinct intervals have an empty gap in which
a cutoff can change from zero to one.

Coordinate lists alone do not determine vector configurations.  To determine
the joint counts, test products of the interval cutoffs just constructed.
There are at most $Q^n$ such tests and their Lipschitz constants are at most
$C/w$.  On the good coordinate fibers every product equals an exact
indicator.  After decreasing the threshold if necessary and discarding a
further relative volume at most $\delta/3$, each
joint count is within less than $1/4$ of its smooth average.  It is an
integer, so it is the same on all remaining fibers.  Taking this many
copies of one point in each occupied coordinate box gives $\mathcal A$.
The boxes have diameter $Cw$, proving the majority assertion.  Intersecting
the good sets of two overlapping balls proves the matching assertion.
Write $d_{\rm sel}(\delta)>0$ for a threshold sufficient for this finite
selection, and use $d_{\rm sel}(1/4)$ when $\delta>1/4$.

There is also an exact finite selection when the energy on this ball is
zero.  Fix $0<\delta'<\min\{\delta,1/2\}$, apply only the finite selection
just proved with positive upper bounds $d_j\downarrow0$ below
$d_{\rm sel}(\delta')$, and use the limsup construction and matching through
one common fiber from the start of this proof.  That part uses only the good sets and their
matching errors, not the moment or support estimates.  It yields an exact
majority configuration with exceptional proportion at most $\delta'<\delta$.

The dyadic covering geometry below, including its enlargement and overlap
ratios, is fixed from the data at the outset and determines $c$ in the
statement.  Choose an internal proportion $0<\delta_{\rm int}<1/4$ so that
two exceptional sets cannot fill any comparison overlap, and small enough
for the support argument.  Then choose a fixed threshold
$d_{\rm int}\le d_{\rm sel}(\delta_{\rm int})$ small enough for that argument
and the iteration, and set
\[
 d_0(\delta_0)=\min\{d_{\rm int},d_{\rm sel}(\delta_0)\}.
\]
For every selection ball in the descent, including the root, use the actual
normalized mass of $\tau$ on the associated fixed enlargement as the selection
upper bound, using the exact finite selection when that mass is zero.
Use the root configuration $\mathcal A$ selected with
$\delta_{\rm int}$ throughout the descent.  If $\delta_0<\delta_{\rm int}$,
apply the finite selection once more at the root with exceptional proportion
$\delta_0$ and the same actual energy.  A common good fiber matches the two
root configurations within
$C_{\delta_0}s\sqrt d$, after enlarging that constant.  On the good fibers
of the second selection, the triangle inequality then gives the prescribed
majority bound for $\mathcal A$.  If $\delta_0\ge\delta_{\rm int}$, the
original good set suffices.  Thus the descent and its root configuration
use only the fixed internal choices; the finer selection affects only the
majority constant and input threshold.

We next record how to include every support point.  On a continuing ball
of radius $r$ with normalized $\tau$ at most $d_{\rm int}$, the majority
assertion and \Cref{lt:scalar,lt:jacobian} bound the mass outside the
configuration's height neighborhoods by
\[
 C(\delta_{\rm int}+d_{\rm int})r^m.
\]
Indeed the mass above bad fibers is at most their $Q$-fold base volume plus
$|M|$, and the projection error also has mass at most $C(D+\Lambda)$.
Because the restricted support is separated from the artificial vertical
boundary, its extension by zero across that boundary remains stationary
for fields supported away from the lateral boundary; multiply such a field
by a vertical cutoff equal to one near the restricted support to verify
this assertion.  The monotonicity ball is taken in this extension and
remains in the lateral working region; it need not lie inside the original
vertical cutoff.  If a support point over an inner base ball were at
distance at least $c_1r$ from these height neighborhoods, its ambient ball
of radius $c_2r$ would lie outside them.  For fixed small
$c_2>0$, this contradicts \eqref{eq:lt-lower-mass} when
$\delta_{\rm int},d_{\rm int}$ are
sufficiently small.  The same reasoning works in the curved tube since its
coordinate map and inverse are uniformly Lipschitz.  This proves an
$O(r)$ bound for every support point on an interior ball.  In particular
it applies to a first stopped child, using its continuing parent.

For the quantitative iteration, subdivide the inner base region into
dyadic cubes, using around each cube a ball of a fixed large multiple of
its side.  Start with sufficiently generous interior room.  Continue down
a branch while the normalized $\tau$ on these balls is at most the fixed
threshold $d_{\rm int}$.  Otherwise stop at the first such cube.  On continuing
balls the normalized mass is bounded by $Q+Cd_{\rm int}$, the integer is $Q$, and
the majority configurations at a child and its parent match with increment
\[
 Cr\sqrt{r^{-m}\tau(\text{fixed enlargement of the parent})}.
\]
The fixed parent enlargement here contains both selection enlargements.
If its $\tau$-mass is zero, both actual selection energies vanish, their
exact majority configurations match on the overlap, and the child cannot stop.
At a first stopped cube the all-support argument gives instead $Cr$ for
the distance from each support height to the parent's configuration.  Its
stopping condition charges $r^m$ to a fixed multiple of $\tau$ on that
enlargement.  Configurations are regarded with their optimal matching, so
their nearest-point distances obey the triangle inequality even at
coincident entries.

Write $r_j=2^{-j}s$, and let $b_j(q)$ be the corresponding increment for a
point above a continuing cube, or its final error at its first stopped
cube; put $b_j=0$ after a stop.  Fixed enlargements of cubes have bounded
overlap at each generation.  The preceding charge, the bounded mass of a
continuing parent, and $\tau(B_{cs})\le d s^m$ give
\begin{align}
 s^{-m}\int b_j^2\,\dd\mu&\le C r_j^2d,
                         \label{eq:lt-increment-two}\\
 \|b_j\|_\infty&\le C r_j
          \min\{1,2^{jm/2}\sqrt d\}.
                         \label{eq:lt-increment-sup}
\end{align}
For the second bound on stopped cubes, stopping itself implies
$2^{jm}d\ge c>0$, so its $Cr_j$ error has the required form.
For a branch which never stops, the all-support $Cr_j$ bound tends to zero;
thus its configurations approach the actual support height in distance.
It follows in both cases that
$\dist(\zeta(q),\mathcal A)\le\sum_jb_j(q)$.

The covering, stopping decisions, and configuration were fixed before
either exponent was chosen.  The following summations therefore apply to
this same $\mathcal A$ for every exponent in the stated ranges.
Interpolate \eqref{eq:lt-increment-two} with the second part of
\eqref{eq:lt-increment-sup}, without the minimum, to obtain
\[
 \left(s^{-m}\int b_j^p\,\dd\mu\right)^{1/p}
 \le Cs\sqrt d\,2^{-j[1-m(1/2-1/p)]}.
\]
Minkowski's inequality and \eqref{eq:lt-p-range} prove
\eqref{eq:lt-config-moment}.  Summing
\eqref{eq:lt-increment-sup} above and below $2^{-j}\simeq d^{1/m}$ gives
$Cs\sqrt d$ when $m=1$, $Cs\sqrt d(1+|\log d|)$ when $m=2$, and
$Cs d^{1/m}$ when $m>2$.  Each is bounded by $C_\beta s d^\beta$ for the
stated range of $\beta$, proving \eqref{eq:lt-config-sup}.
\end{proof}

\subsection{Gaussian and height identities}

\begin{corollary}[Gaussian weights and polynomial dilations]
\label{lt:weightedballs}
Define
\[
 G_s(y)=(2\pi s^2)^{-m/2}\exp(-|y|^2/(2s^2))
\]
and fix a center $x$.
On the region $|y-x|\le Rs$, cover the base by balls of radii comparable to
$s/(1+|y-x|/s)$, with fixed enlarged balls having bounded overlap.
Suppose each enlarged ball meets the hypotheses of \Cref{lt:scalar} with
the same propagated integer $Q$.  Then
\begin{equation}\label{eq:lt-weighted-scalar}
 \int_{|y-x|\le Rs}G_s(y-x)\,\dd|M|(y)
 \le C\int_{|y-x|\le (R+C)s}G_s(y-x)\,\dd(D+\Lambda)(y).
\end{equation}
The constant is independent of $R$.  The part outside the indicated region
is retained as a tail error.

For the configuration conclusion, fix one exceptional proportion
$\delta_{\rm cmp}>0$ smaller than the fixed comparison-overlap proportions.
If, in addition, on all admissible dilations $1\le t\le R_*$ one has
\begin{equation}\label{eq:lt-polynomial-dilation}
 (ts)^{-m}\tau(B_{cts}(x))\le C t^Jd,
\end{equation}
with $Ct^Jd\le d_0(\delta_{\rm cmp})$ throughout this range and the same
integer propagated, then configurations
matched from the scale $s$ have moment and support estimates on a dilation
$t$ equal to those in \Cref{lt:config}, with an additional fixed power of
$t$.  Gaussian integrals of these estimates have constants depending on
$J$ and the fixed moment order, once the tail beyond $R_*$ is accounted for.
\end{corollary}

\begin{proof}
The logarithmic derivative of $G_s$ is $-(y-x)/s^2$.  Its values on each
fixed enlargement of a ball of the stated radius differ by bounded factors.
Such a cover can be obtained by a maximal disjoint family of sufficiently
small balls for this radius function; nearby radii are comparable.  Apply
\eqref{eq:lt-scalar-mass} on each ball and sum.  This proves
\eqref{eq:lt-weighted-scalar} without a factor depending on the number of
balls.  Apply \Cref{lt:config} with the fixed proportion $\delta_{\rm cmp}$
successively to concentric dyadic dilations
and to overlapping comparable balls.  Every increment is bounded by a
fixed power of $t$ times $s\sqrt d$, and each local moment estimate has the
same property by \eqref{eq:lt-polynomial-dilation}.  Summation of the
geometric increments gives the claimed polynomial bounds.  Finally
$\sum_{j\ge0}2^{jC}\exp(-c4^j)<\infty$ absorbs any fixed polynomial loss
in the Gaussian integral.
\end{proof}

We specify the cutoff convention for every Gaussian identity.  Choose a
smooth base cutoff $\chi$ which is one on the region under consideration
and has compact support in the chart.  A subscript $s$ denotes convolution
at $x$ with $G_s$ of the extended measure; on the chart it is the original
measure, and outside it any stated extension is used.  In a first-variation
identity the test field is multiplied by $\chi$.  Contributions where
$\chi\ne1$, including the difference from the chosen extension, are denoted
by $\mathcal T_s$.  If the cutoff annulus has distance at least $h$ from $x$
and the relevant total measure is at most $K_0$, decrease $h$ to the
available transition width if necessary, and choose $0\le\chi\le1$ with
$\|\nabla\chi\|_\infty\le C(1+h^{-1})$.  The first-variation product rule
differentiates this fixed cutoff once; scale derivatives leave it fixed.
For $0<s\le1$, direct differentiation of the Gaussian bounds each such term by
\begin{equation}\label{eq:lt-chart-tail}
 C K_0 s^{-C}(1+h^{-1})^C\exp(-h^2/(4s^2)),
\end{equation}
times the stated bounds for the height multiplier and its derivatives.
One may use a larger fixed exponent $C$ for the finite set of Gaussian and
height-multiplier derivatives used below.
When a polynomial dilation bound such as
\eqref{eq:lt-polynomial-dilation} is available through $R_*$, annular
summation bounds the contribution from $Rs<|y-x|\le R_*s$ by
$C\,\mathrm{poly}(R)e^{-cR^2}$ times its reference density.  The part
beyond $R_*s$ remains a separate tail, including any part still inside
the chart, until a further dilation or total-mass bound controls it.
The exterior-chart term is also retained.  A periodized
Gaussian on a fixed large torus obeys the same estimates.  These formulas
are the meaning of a negligible Gaussian tail; no unrestricted extension
is silently discarded.

\begin{lemma}[Gaussian mass identity]\label{lt:gaussian}
Set $u=(y-x)/s$, $A(s)=s^2M_s(x)$, and
\[
 R_s(x)=\int G_s(y-x)B[u,u]\,\dd\mu.
\]
In flat coordinates, apart from $\mathcal T_s$,
\begin{equation}\label{eq:lt-gaussian-flat}
 \dot A=s^2L_s+s^2R_s,
 \qquad \dot{}=\frac{\dd}{\dd\log s}.
\end{equation}
In the curved coordinates fixed above,
\begin{equation}\label{eq:lt-gaussian-curved}
 \dot A=s^2L_s+s^2R_s+\mathcal E_s+\mathcal T_s,
 \qquad
 |\mathcal E_s|\le C\gamma s^2
           \int G_s(y-x)(1+|u|^2)\,\dd\tau(y).
\end{equation}
Here the entries of $u$ in $B[u,u]$ are taken in the chosen orthonormal
horizontal frame, and $\gamma$ retains its definition from the fixed
physical coordinates.
\end{lemma}

\begin{proof}
In flat coordinates $\dot G_s=-\Div_y((y-x)G_s)$.  The horizontal first
variation says $\int\Div X\,\dd M=\int B:DX\,\dd\mu$.
Take $X=(y-x)G_s$.  Since
$DX=G_s(I-u\otimes u)$, it follows that
$\dot M_s=-D_s+R_s$.  Differentiating the factor $s^2$ proves
\eqref{eq:lt-gaussian-flat}.

For the curved identity apply \Cref{lt:curvedvariation} to the same
coordinate field.  Its scaled derivative is bounded by
$CG_s(1+|u|^2)$ and its size by $CsG_s|u|$.
On a physical chart of radius $l$, coordinate coefficients and their
scaled first derivatives differ from the centered Euclidean ones by
$O(l)$.  The divergence comparison is made after subtracting
$Q\vol_S$; its coefficient error therefore multiplies $M$, not the
background mass.  The tensor coefficient errors multiply $B$ and are
bounded by $C\gamma e$.  Equation \eqref{eq:lt-scaledvariation} controls
the remaining geometric terms.  Multiplying by $s^2$ proves
\eqref{eq:lt-gaussian-curved}.  Differentiation of $\chi$ produces precisely
the tails described in \eqref{eq:lt-chart-tail}.
\end{proof}

\begin{lemma}[Height tests and geometric errors]\label{lt:heighttests}
In flat coordinates, for a base test $\varphi$, a scalar height function
$f$, a horizontal field $X$, and a normal vector function $v$, stationarity
gives
\begin{align}
 \int f\Div X\,\dd\mu
 &=\int f B:DX\,\dd\mu-
                   \int Df\cdot P_{vh}X\,\dd\mu,
                     \label{eq:lt-height-horizontal}\\
 0&=\int\{\varphi Dv:P_{vv}
                +v\cdot P_{vh}\nabla\varphi\}\,\dd\mu.
                     \label{eq:lt-height-normal}
\end{align}
All height functions here use the scaled variable $\zeta$.
For a scale-independent $C^2$ height function $\Gamma$, these formulas give
\begin{equation}\label{eq:lt-height-gaussian}
 \frac{\dd}{\dd\log s}\int G_s\Gamma\,\dd\mu
 =\int G_s\Gamma(B[u,u]-e)\,\dd\mu
                       +s^2\int G_sD^2\Gamma:P_{vv}\,\dd\mu,
\end{equation}
apart from the specified chart tails.  If $\Gamma$ depends on $s$, add
$\int G_s\dot\Gamma\,\dd\mu$.

In the curved tube, use the original radial-gauge normal frame fixed above.
The extra trace terms in a normal test
$\varphi v(\zeta)$ without a derivative of $\varphi$ are bounded pointwise
by $C|\varphi|$ times
\begin{equation}\label{eq:lt-normal-error}
 |Dv|l|z_{\rm phys}|\sqrt e
       +l|v|\bigl(|H_S|+|z_{\rm phys}|+e+\sqrt e\bigr).
\end{equation}
The terms with a derivative of $\varphi$ have the flat tensor expression
in \eqref{eq:lt-height-normal}, with a coefficient error at most
$C\gamma|P_{vh}|$.
The curved horizontal identity is covariant, with $\Div_SX$ and
$\nabla^SX$ in place of the flat base operators.  Its geometric error is
bounded by the right side of \eqref{eq:lt-scaledvariation} with the integrand
multiplied by $|f|$.  The additional flux error from the height gradient is
bounded pointwise by
\begin{equation}\label{eq:lt-horizontal-height-error}
 C\gamma |Df||X|(\sqrt e+\sqrt\lambda).
\end{equation}
In its coordinate form the divergence comparison also contributes an error
bounded by $Cl\int|f||X|\,\dd\mu$; the coordinate changes in the $B$ term
are included in the preceding geometric error.
Under a further interior rescaling, retain the original horizontal and
normal frames and the original $\gamma$ in \eqref{eq:lt-small-geometry}.
Transform the displayed estimates by change of variables.  The inherited physical
connection remains $O(l)$; no improvement to $O(ls)$ is asserted on an
off-center ball of coordinate radius $s$.  In particular
first-order flux errors have size at most
$C\gamma(\sqrt e+\sqrt\lambda)$ times the relevant multiplier, whereas
quadratic errors have size at most $C\gamma(e+\lambda)$.  The factor
$|v|$ in \eqref{eq:lt-normal-error} must be retained when estimating a
normal test.
\end{lemma}

\begin{proof}
In a product Euclidean chart, differentiate the ambient fields
$f(\zeta)X(y)$ and $\varphi(y)v(\zeta)$.  Their horizontal and vertical
blocks give \eqref{eq:lt-height-horizontal} and
\eqref{eq:lt-height-normal}.  Apply the first with
$f=\Gamma$ and $X=(y-x)G_s$.  It yields the first term of
\eqref{eq:lt-height-gaussian} and the flux
$s\int G_sD\Gamma\cdot P_{vh}u\,\dd\mu$.
Apply the normal formula with $v=D\Gamma$ and $\varphi=G_s$; since
$\nabla G_s=-G_su/s$, this flux is exactly
$s^2\int G_sD^2\Gamma:P_{vv}\,\dd\mu$.  This proves the identity and
fixes both its sign and its factors of $s$.

We give the geometric error calculation.  In radial gauge the normal
connection one-form vanishes at the chart center and has physical norm
$O(l)$ throughout the chart, by its bounded curvature and the radial
parallel-transport equation.  Its ambient derivative still has the
tangential component $-A_v$; that component is not part of the small
connection one-form.  A horizontal derivative in scaled ambient
coordinates of $\zeta=z_{\rm phys}/l$ is
\[
 D_h\zeta=-\omega(J_z^{-1}h)z_{\rm phys},
 \qquad |D_h\zeta|\le Cl|z_{\rm phys}||h|,
\]
while the vertical derivative is unchanged.  Differentiating $v(\zeta)$
therefore gives the first summand of \eqref{eq:lt-normal-error} on
contraction with the mixed block.  Differentiating its normal frame gives
a horizontal trace whose identity part is
$-l\langle H_S,v\rangle$, with error $Cl|z_{\rm phys}||v|$ from
$J_z^{-1}-I$; its $B$ part is bounded by $Cle|v|$.
The remaining mixed-frame term is bounded by $Cl\sqrt e|v|$.
This proves \eqref{eq:lt-normal-error}.  A base derivative uses
$D\pi=J_z^{-1}P_h$; its coefficient on a coordinate gradient in the chosen
frame is $J_z^{-1}\mathsf C_y^{-\top}$, which differs from $I$ by at most
$C\gamma$.  This gives the stated coefficient error on the mixed block.

For the horizontal test, first use the calculation in
\Cref{lt:curvedvariation}, and then differentiate $f(\zeta)$.
If $X_{\rm coord}$ is the coordinate column of the scaled field, its lift
has column $J_z\mathsf C_yX_{\rm coord}$ in the horizontal frame.  Subtracting
the flat mixed flux leaves
\[
 Df\cdot P_{vh}(J_z\mathsf C_y-I)X_{\rm coord},
 \qquad |J_z\mathsf C_y-I|\le C\gamma.
\]
The bound $|P_{vh}|\le C\sqrt e$ gives its first contribution to
\eqref{eq:lt-horizontal-height-error}.  The separate connection derivative
is bounded by $Cl|z_{\rm phys}||Df||X_{\rm coord}|$, which gives the second
contribution because $l\le\gamma$ and $|z_{\rm phys}|\le\sqrt\lambda$.
For the coordinate form, in the original centered coordinates put
$a_l(\xi)=\sqrt{\det g(l\xi)}$.  For a scaled coordinate field $X(\xi)$,
\[
 \Div_S(l\,DF(l\xi)X)=\Div_\xi X+D_\xi(\log a_l)\cdot X,
 \qquad |D_\xi\log a_l|\le Cl.
\]
This weighted identity has no subtracted reference volume, so this
last term gives the stated $Cl\int|f||X|\,\dd\mu$ error.  The remaining
coordinate coefficient changes are contracted with $B$ and cost
$Cl e|f|(|\nabla X|+|X|)$, already included in the geometric error.
At a further radius $s$, put $y'=(y-x)/s$, $\zeta'=\zeta/s$,
$T_{x,s}(y,\zeta)=(y',\zeta')$, $\mu'=s^{-m}(T_{x,s})_\#\mu$, and
$v'(\zeta')=v(s\zeta')/s$.  Then $D'v'=Dv$, while $\lambda$ remains in
physical units.  Keeping the same frame retains the coefficient $l$ in
$l|Dv||z_{\rm phys}|\sqrt e$; the terms involving $v$ become
$(ls)|v'|(|H_S|+|z_{\rm phys}|+e+\sqrt e)$.
The resulting estimates use the original $\gamma$.  All configuration
comparisons likewise use the original height frame, so there is no change
of height coordinates from a new gauge.
Young's inequality changes products of two first-order errors into the
stated quadratic bound.  Applying Cauchy--Schwarz to
\eqref{eq:lt-normal-error} retains the actual $L^2$ size of $v$, which is
why this form of the bound is useful for displacement and clipped-height
tests.  Finally multiply each field by the chart cutoff; its derivative
terms obey \eqref{eq:lt-chart-tail}.
\end{proof}
